\documentclass[9.5pt,a4paper,landscape]{article}
\usepackage[utf8]{inputenc}
\usepackage[T1]{fontenc}
\usepackage[french]{babel}
\usepackage[a4paper,landscape,margin=0.8cm]{geometry}
\usepackage{lmodern}
\usepackage{xcolor}
\usepackage{tcolorbox}
\usepackage{array}
\usepackage{amsmath}

\pagestyle{empty}
\setlength{\parindent}{0pt}

\definecolor{bleu}{HTML}{245B78}
\definecolor{gris}{HTML}{F7F7F7}

\newtcolorbox{exo}[1]{
  colback=white,
  colframe=black!55,
  boxrule=0.45pt,
  arc=1.2mm,
  left=1.8mm,
  right=1.8mm,
  top=1.5mm,
  bottom=1.3mm,
  height=5.35cm,
  valign=top,
  title=\textbf{Exercice hors classe #1},
  coltitle=bleu,
  fonttitle=\small,
  colbacktitle=gris,
  titlerule=0.25pt
}

\begin{document}

\fontsize{9}{10.5}\selectfont

\begin{center}
{\large\bfseries\color{bleu}
Séquence 1 -- Statistiques : exercices hors classe}\par

\vspace{0.1em}

{\small
À faire dans le cahier.
Chaque exercice est court et reprend directement un travail déjà rencontré en classe.}
\end{center}

\vspace{0.25em}

%==================================================
% LIGNE 1
%==================================================

\noindent
\begin{minipage}[t]{0.242\linewidth}
\begin{exo}{1}
Dans un parking de 50 voitures, on compte 8 voitures blanches,
12 noires et 14 bleues. Les autres sont rouges.

\medskip
Combien y a-t-il de voitures rouges ?
\end{exo}
\end{minipage}
\hfill
%
\begin{minipage}[t]{0.242\linewidth}
\begin{exo}{2}
Dans une classe de 24 élèves, 6 viennent au collège à vélo.

\medskip
Donne la fréquence correspondante sous forme de fraction puis de pourcentage.
\end{exo}
\end{minipage}
\hfill
%
\begin{minipage}[t]{0.242\linewidth}
\begin{exo}{3}
Dans un groupe de 32 élèves, 25\% ont choisi l'option théâtre.

\medskip
Combien d'élèves cela représente-t-il ?
\end{exo}
\end{minipage}
\hfill
%
\begin{minipage}[t]{0.242\linewidth}
\begin{exo}{4}
Voici cinq nombres :

\[
8\quad 11\quad 9\quad 12\quad 10
\]

Calcule leur moyenne.
\end{exo}
\end{minipage}

\vspace{0.16cm}

%==================================================
% LIGNE 2
%==================================================

\noindent
\begin{minipage}[t]{0.242\linewidth}
\begin{exo}{5}
Une joueuse a obtenu :

\begin{center}
\begin{tabular}{c|ccc}
Score & 8 & 10 & 12\\
\hline
Effectif & 2 & 3 & 1
\end{tabular}
\end{center}

Calcule son score moyen.
Arrondis au dixième si besoin.
\end{exo}
\end{minipage}
\hfill
%
\begin{minipage}[t]{0.242\linewidth}
\begin{exo}{6}
Détermine la médiane de la série :

\[
11;\ 17;\ 14;\ 26;\ 19;\ 30;\ 12
\]
\end{exo}
\end{minipage}
\hfill
%
\begin{minipage}[t]{0.242\linewidth}
\begin{exo}{7}
Détermine la médiane de la série :

\[
13;\ 21;\ 19;\ 18;\ 27;\ 15
\]
\end{exo}
\end{minipage}
\hfill
%
\begin{minipage}[t]{0.242\linewidth}
\begin{exo}{8}
La médiane d'une série est 14.

\medskip
Complète :

Au moins la moitié des valeurs sont \dotfill{} 14
et au moins la moitié sont \dotfill{} 14.
\end{exo}
\end{minipage}

\vspace{0.16cm}

%==================================================
% LIGNE 3
%==================================================

\noindent
\begin{minipage}[t]{0.242\linewidth}
\begin{exo}{9}
Voici des profondeurs atteintes en apnée :

\[
-82;\ -96;\ -86;\ -92;\ -87\ \text{m}
\]

Quelle est la profondeur médiane ?
\end{exo}
\end{minipage}
\hfill
%
\begin{minipage}[t]{0.242\linewidth}
\begin{exo}{10}
On interroge 28 élèves sur le nombre de frères et sœurs qu'ils ont.

\medskip
Indique :
\begin{itemize}
  \item la population étudiée ;
  \item le caractère étudié.
\end{itemize}
\end{exo}
\end{minipage}
\hfill
%
\begin{minipage}[t]{0.242\linewidth}
\begin{exo}{11}
Dans un groupe de 40 élèves, 14 viennent au collège en bus.

\medskip
Donne la fréquence correspondante :
\begin{itemize}
  \item sous forme de fraction ;
  \item en écriture décimale ;
  \item en pourcentage.
\end{itemize}
\end{exo}
\end{minipage}
\hfill
%
\begin{minipage}[t]{0.242\linewidth}
\begin{exo}{12}
La durée moyenne du trajet domicile--collège d'un groupe d'élèves
est de 18 minutes.

\medskip
Écris une phrase qui interprète cette moyenne.
\end{exo}
\end{minipage}

%==================================================
% CORRIGÉ PROF (non imprimé)
%==================================================
%
% 1) 16.
% 2) 6/24 = 1/4 = 25 %.
% 3) 8.
% 4) 10.
% 5) (2x8 + 3x10 + 1x12)/6 = 58/6 ≈ 9,7.
% 6) 17.
% 7) 18,5.
% 8) inférieures ou égales à / supérieures ou égales à.
% 9) -87 m.
% 10) Population : les 28 élèves interrogés.
%     Caractère : le nombre de frères et sœurs.
% 11) 14/40 = 7/20 = 0,35 = 35 %.
% 12) Interprétation possible :
%     si la durée totale des trajets était répartie également
%     entre les élèves, chacun aurait un trajet de 18 minutes.
%
\end{document}
